\documentclass[10pt,a4paper]{article}
\usepackage[T1]{fontenc}
\usepackage[utf8]{inputenc}
\usepackage{lmodern}
\usepackage{amsmath,amssymb}
\usepackage[margin=21mm]{geometry}
\usepackage{microtype}
\usepackage[hidelinks]{hyperref}
\usepackage{enumitem}
\setlist[itemize]{leftmargin=*,nosep,topsep=3pt}
\setlength{\parindent}{0pt}
\setlength{\parskip}{5pt}
\setlength{\emergencystretch}{3em}
\newcommand{\question}[1]{\par\medskip{\normalsize\bfseries #1}\par\smallskip}
\pagestyle{plain}

\begin{document}
{\Large\bfseries Generalized sleep decoding with basal ganglia signals in multiple movement disorders\par}

\textbf{Paper:} Yin et al. (2024).\quad Human sleep-decoding study.\par
\href{https://doi.org/10.1038/s41746-024-01115-7}{DOI: \nolinkurl{10.1038/s41746-024-01115-7}}\par
{\footnotesize Page references below use PDF page numbers, including front matter.}

\question{To which general question does this paper relate?}
What implanted electrodes can measure, and how feedback could account for sleep. Evaluates a sleep-state estimator across people, diagnoses and recording targets.

\question{How does it answer the general questions?}
\textbf{Measured.} Synchronized polysomnography and intracranial field potentials provided training labels and neural features. After exclusions, 140 nights from 121 patients were analysed. BGOOSE achieved approximately 85\% average accuracy for awake/NREM/REM classification; additional ECoG improved decoding (PDF pp.~1--7).

\textbf{Usable insight.} Cross-subject evaluation and external datasets support using neural sensing to estimate context without individual polysomnography calibration. Performance is not uniform across stages: separating NREM substages is substantially harder. Classification results do not themselves show that sleep-triggered stimulation improves sleep (PDF pp.~6--8, 11).

\textbf{Model relevance.} A sleep estimator could select operating conditions or feedback policies. Its spectral, waveform and entropy features require an observation/processing layer beyond a linear $y=Cx$. No stimulation-to-neural-state dynamics or delayed-model coefficients are identified. Sleep quality would require its own validated performance output.

\question{What papers does it cite that could provide more information?}
\begin{itemize}
\item Smyth et al. (2023), \emph{Adaptive deep brain stimulation for sleep stage targeting in Parkinson's disease}: intervention-oriented follow-up (ref.~12, PDF p.~12).
\item Gilron et al. (2021), \emph{Long-term wireless streaming of neural recordings for circuit discovery and adaptive stimulation in individuals with Parkinson's disease}: chronic sensing (ref.~28, PDF p.~12).
\end{itemize}

\end{document}
